\bookchapter{第二十一章：モバイルネットワークと情報セキュリティの深掘り}

\chapterintro{通常のネットアクセス}

\booksection{：あなたのスマホ公共ネットワーク対象サイト。中間のルーター、スイッチ、どの段階でもあなたが送信するデータを「覗き見」される可能性。特にカフェ、空港のような公共Wi-Fiでは、正直「裸走」状態。}

VPN使用

\begin{itemize}
  \item \textbf{：あなたのスマホVPN暗号化トンネルVPNサーバー対象サイト。データは暗号化カプセル化され、「覗き見」者は乱数しか見えない。また、対象サイトが見るのはVPNサーバーのIPアドレスで、あなたの本当のIPではありません。}VPNで何ができるか？簡単に言うと：第一、
  \item \textbf{公共Wi-Fiでのプライバシー保護}。第二、
\end{itemize}

地域制限の回避\textbf{（特定の国でのみ配信されているコンテンツの視聴など）。第三、}ネットワークプロバイダー（ISP）の閲覧習慣記録の防止\textbf{。ただし一言注意：}無料VPNを使わないでください！\textbf{無料の背後はしばしばあなたのデータが彼らの「商品」であることを意味します。ある友人が無料VPNを使ったら、スマホに謎の広告がポップアップし始め、後にそのVPNが彼の閲覧データを広告業者に売却していたことが判明。だからVPN選びは、有料の大手ブランド（NordVPN、ExpressVPNなど）を選ぶか、使わないか。}また、VPNは万能ではありません。「転送中」のデータしか保護できず、対象サイト自体が安全でない、またはスマホにマルウェアが感染している場合、VPNは助けになりません。\textbf{二、公共Wi-Fiの「甘い罠」——想像より危険}無料Wi-Fiに接続するだけだろ、どれほどのリスクがあると思う？実話をお話しします：セキュリティ研究者が実験を行い、空港に「Starbucks\_Free」というWi-Fiホットスポットを設置、1時間足らずで百人以上が接続。実際にはこのホットスポットは研究者が自分で設置したもので、このホットスポットを経由するすべての

暗号化されていないトラフィック

\booksection{——例えばあなたがアクセスしたHTTPウェブページの内容、TLS未対応のApp通信——は彼らが傍受·閲覧できました。}

ただし重要な前提：今日日常的に使っているWeChat、ネットバンキング、Alipayは、ほとんどがHTTPS暗号化を採用。つまり、攻撃者がデータパケットを傍受しても、見えるのは乱数だけで、\textbf{チャット内容や決済パスワードを直接読み取ることはできない}。ただし彼らはあなたがどのサイトにアクセスしたか（ドメイン名）は依然として分かります。一部の古いAppやウェブページがまだHTTP平文送信を使っている場合、内容は完全に露出。この攻撃を「中間者攻撃」（Man-in-the-Middle）と呼び、原理は攻撃者があなたと本当のWi-Fiルーターの間に「割り込む」こと。さらに恐ろしいのは、多くのスマホが「既知Wi-Fiに自動接続」を設定しており、攻撃者が自分のホットスポット名をあなたが以前接続したWi-Fi（例えば「ChinaNet」や「CMCC」）に設定するだけで、あなたのスマホが自動的に接続。

どう防ぐか？いくつかの実用的なアドバイス：\textbf{HTTPSを確認}：アドレスバーの先頭がhttps://で、余分な「s」はSecureを意味。パスワード入力が必要なサイトにアクセスする時、http://（sなし）を見たら、すぐに止まって入力しない！最近の主流ブラウザは非HTTPSサイトに「安全ではない」警告を表示、このマークを見たら警戒を。

「自動接続」をオフ

\begin{itemize}
  \item \textbf{：iPhoneでもAndroidでも、Wi-Fi設定で「自動参加」または「自動接続」をオフに。毎回手動でタップする方が、スマホに判断させるより良い。}スマホのモバイルデータで機密操作
  \item \textbf{：送金、ネットバンキングログイン、ID番号入力……これらの操作は自分の4G/5Gモバイルデータで着実に行い、あの通信量を節約しようとしない。}使わないWi-Fiを忘れる
  \item \textbf{：出張や旅行から戻ったら、以前一時的に接続したホテルWi-Fi、空港Wi-Fiをすべて「忘れる」。これで次回スマホが同名の偽ホットスポットに自動接続しない。}三、HTTPSはどうあなたを保護する？——SSL/TLSハンドシェイク
  \item \textbf{先ほどHTTPSに言及しましたが、実際にどう動作するのでしょう？HTTPSサイトにアクセスするたび、ブラウザとサーバー間で「SSL/TLSハンドシェイク」が行われ、全体は大体こうです：}ブラウザが「ねえ、サーバー、安全な通信したいんだけど、これが私の暗号化能力リスト（対応暗号化アルゴリズムなど）」と言う。
\end{itemize}

\booksection{サーバーが返答：「OK、あなたのリストからこの暗号化方案を選んだ、これが私のデジタル証明書、公開鍵が含まれ、某某CA機関が署名認証」。}

ブラウザが証明書の真偽を検証：信頼できるCAが発行？期限切れ？ドメインが一致？

\begin{enumerate}
  \item 検証通過後、双方が公開鍵暗号で「セッションキー」をネゴシエート、以降の通信はすべてこのキーで暗号化。
  \item これは銀行で手続きするようなもの：まず身分証（証明書）を提出、窓口があなたの身分を確認、それから二人で合言葉を決め、以降の会話は二人だけが理解できる。誰かが傍聴していても、暗号化された乱数しか聞こえない。
  \item ただし注意：HTTPSは「転送過程」が暗号化されていることだけを保証し、対象サイトが善人であることは保証しない。フィッシングサイトもHTTPSを使える——あの鍵アイコンは「あるサイト」に接続していることだけを示し、そのサイトがあなたが行きたいサイトであることは保証しない。だから自分でURLが正しいか確認する必要があります。
  \item 四、DNSセキュリティ：あなたの「ネットワーク電話帳」は安全？
\end{enumerate}

DNS（ドメイン名システム）はインターネットの電話帳のようなもの：あなたが「baidu.com」と入力すると、DNSが対応するIPアドレス（例えば110.242.68.66）を調べ、ブラウザが接続できます。問題は、伝統的なDNSクエリは

平文

\booksection{で、}

暗号化されていない\textbf{！つまり、あなたのISP、Wi-Fi提供者、または中間ノードは誰でもあなたがどのサイトにアクセスしたか見える。}さらに厄介なのは「DNSハイジャック」：攻撃者がDNS段階で細工し、「baidu.com」を偽のフィッシングサイトに向ける。あなたが正しいURLを入力しても、全く同じだが実際は偽のログインページに連れて行かれ、アカウントパスワードを入力した瞬間に終わり。\textbf{幸い暗号化DNS方案が登場：}DNS over HTTPS（DoH）

：DNSクエリをHTTPSトラフィックに隠し、通常のウェブ閲覧トラフィックに混ぜる、どれがDNSクエリか区別できない。

DNS over TLS（DoT）

\begin{itemize}
  \item \textbf{：TLS暗号化プロトコルでDNSクエリを専用保護、ポートが独立、ネット管理者はDoT使用が分かるが内容は見えない。}私はCloudflareの1.1.1.1またはGoogleの8.8.8.8をDNSサーバーとして推奨、両方とも暗号化DNSに対応。Androidの「プライベートDNS」設定で「one.one.one.one」または「dns.google」を直接入力するだけで、超簡単。iPhoneは少し手間がかかり、プロファイルのインストールまたはApp設定が必要だが、やる価値あり。
  \item \textbf{五、「シークレットブラウジング」は本当に透明になれる？}まず結論：
\end{itemize}

シークレットモードはあなたの家のカーテンであって、透明マントではありません。

\booksection{シークレットモードがすること：ローカルに痕跡を残さない——シークレットウィンドウを閉じると、ブラウザは閲覧履歴、Cookie、フォーム入力記録、検索記録を保存しない。他人のパソコンで何かを調べる、プレゼントを買って家族に閲覧履歴から見つかりたくない場合、シークレットモードが適切。}

折りたたみ画面は「未来のテクノロジー」から「今買える」成熟カテゴリーになりました。しかし、ストレート型スマートフォンとはフォームファクター、使用習慣、メンテナンス方法に大きな違いがあり、買う価値があるか、どう選ぶか、どう使うか、この記事で整理しましょう。\textbf{あなたがアクセスしたサイトは}

あなたが来たことを同じように見える

、あなたのIPアドレス、デバイス情報、閲覧行動はサイト側で完全に把握。

\begin{itemize}
  \item あなたのネットワークプロバイダーは\textbf{あなたがどのサイトにアクセスしたか同じように見える}、ただし具体的な内容は見えない（HTTPS使用場合）。
  \item 会社/学校のネットワーク管理者は\textbf{同じように監視できる}あなたのネットアクセス行動。
  \item あなたがダウンロードしたファイル、追加したブックマークは\textbf{残る}、ウィンドウを閉じても消えない。
  \item だから覚えて：シークレットモードは「ローカル消去器」で、「ネットワーク透明着」ではない。本当にネットアクセス痕跡を隠す必要があるなら、VPN＋Torブラウザのようなコンボが必要ですが、日常使用ではこれらを理解していれば十分。\textbf{六、ブラウザプライバシー：ブラウザにあなたを売らせない}気づいていないかもしれませんが、ブラウザ自体が巨大な情報収集器。Cookie、閲覧履歴、キャッシュ、自動入力、位置権限……それぞれがあなたのプライバシーを漏洩する可能性。いくつかのアドバイス：
\end{itemize}

定期的に閲覧データをクリア

\booksection{：定期的に設定でCookie、キャッシュ、履歴をクリア。「追跡防止」のためではありません——Cookieクリアではもうあまり効かない——スマホが重くなるのを防ぐため。ブラウザキャッシュは数GBの容量を占めることがある。}

「Do Not Track」の幻想を捨てる

\begin{itemize}
  \item \textbf{：ブラウザの「Do Not Trackリクエストを送信」オプションは、本質的にサイトに「追跡しないで」とお願いするもの。しかしサイトが聞くかどうかは全く別の話で、ほとんどのサイトは無視。期待しないで。}プライバシー保護ブラウザを活用
  \item \textbf{：Firefox Focus、DuckDuckGoブラウザ、Braveはデフォルトでトラッカーと広告をブロック、ChromeやSafariのデフォルト設定よりずっとクリーン。}検索エンジンも重要
  \item \textbf{：Baidu/Googleの代わりにDuckDuckGoで検索、検索履歴を記録せず、ユーザープロファイルを構築しない。正直検索結果はGoogleほど精確ではないが、日常調べ物には十分。}App権限を確認
  \item \textbf{：スマホ設定でブラウザに付与された権限を確認。本当に位置、マイク、連絡先へのアクセスが必要？おそらく不要。}七、実用的なセキュリティチェックリスト：今から自分を守る
  \item \textbf{多くを語りましたが、最後にシンプルなセキュリティチェックリストを提供します。30分で完了すれば、スマホのセキュリティは質的に向上：}Wi-Fi自動接続をオフ
\end{itemize}

\booksection{：iPhoneは設定-Wi-Fiで、AndroidはWi-Fi設定で。}

暗号化DNSに変更

\begin{enumerate}
  \item \textbf{：AndroidはプライベートDNSを「dns.google」に直接設定、iPhoneはAppまたはプロファイルで1.1.1.1のDoHを設定。}二要素認証（2FA）を有効化
  \item \textbf{：WeChat、Alipay、Apple ID、Googleアカウント、2FAを有効にできるものはすべて有効化。SMS認証コードは最低レベルの2FAで、TOTP（Google Authenticator、Microsoft Authenticatorなど）が使えればさらに良い。}App権限を一通り確認
  \item \textbf{：各Appの権限を確認、特に位置、連絡先、SMS、マイク、カメラの5項目。懐中電灯Appが連絡先を要求するようなものを見つけたら——すぐアンインストール。}定期的にブラウザデータをクリア
  \item \textbf{：カレンダーリマインダーを設定、毎月1回、キャッシュ、Cookieをクリア。}公共充電スタンドのケーブルを使わない
  \item \textbf{：USBポートは充電もデータ転送もでき、悪意のある充電スタンド（「Juice Jacking」と呼ばれる）が充電中にデータを盗む可能性。外出時は自分の充電ヘッドとケーブルを持参。}今すぐやる：スマホに「セキュリティ健康診断」を実施
  \item \textbf{DNSアップグレード}：Android：設定 > 接続と共有 > プライベートDNS > 「one.one.one.one」または「dns.google」を入力。iPhone：Safariで「1.1.1.1 DNS profile」を検索しプロファイルをインストール、またはCloudflareの1.1.1.1 Appをダウンロード。
\end{enumerate}

\begin{tcbpractice}{Wi-Fi監査}
\begin{enumerate}
  \item \textbf{：Wi-Fiリストを開き、知らない、使わない、名前の怪しいホットスポットをすべて「忘れる」。}権限大掃除
  \item \textbf{：設定 > プライバシー/権限管理に入り、Appごとに権限を審査。権限要求が不合理（電卓が連絡先を要求など）なAppをマークし、アンインストールまたは代替品を検討。}アカウントセキュリティを確認
  \item \textbf{：Apple IDまたはGoogleアカウントのセキュリティページを開き、不明なデバイスのログインがないか確認、2FAが有効か確認、復元連絡先が信頼できるか確認。}第二十二章：AIはスマホでどう動くのか
  \item \textbf{スマホの最もクールな「魔術ショー」へようこそ、主役は：人工知能。あなたが毎日使っている「ワンタップで綺麗な写真」、口を開けば応える音声アシスタント、使うほどにあなたを理解する入力メソッド、その背後にこの見えない頭脳が立っています。}一、AIカメラ：あなたの「現像室」がスマホに
\end{enumerate}

\end{tcbpractice}
